\section{Process control and checkpoints}

\subsection{\code{COMPILEPROCESS} (\implemented)}

\begin{lstlisting}
COMPILEPROCESS ChildName
\end{lstlisting}
COMPILEPROCESS looks up the named process and compiles it on a scratch state.
Its gates are not inserted into the current circuit. Use RUNCHILD or CALL when
the child should actually run. An unknown name is an error. A process that
asks to compile itself is verified once and does not recurse.

\subsection{\code{MASTERVAL} (\implemented)}

\begin{lstlisting}
MASTERVAL Result
\end{lstlisting}
When the process has no \code{RETURNVALS}, MASTERVAL names are the ordered
public outputs. When \code{RETURNVALS} is present, every MASTERVAL name must
already appear there, and the truth-table columns stay the RETURNVALS list.

\subsection{Save-state checkpoint instruction (\implemented)}

\begin{lstlisting}
SAVE_STATE checkpointA
\end{lstlisting}
\code{SAVE_STATE} copies the current amplitude vector and measurement record
into a classical store under that name. It is not a unitary and does not
clone the quantum state inside the register.

\subsection{Load-state checkpoint instruction (\implemented)}

\begin{lstlisting}
LOAD_STATE checkpointA
\end{lstlisting}
\code{LOAD_STATE} replaces the active state vector and measurement record with
the named checkpoint. A missing name is an error. Later gates continue from
the restored amplitudes.

On the canvas neither instruction is a gate box. Each is a dashed vertical
line across every wire and the \textbf{c} lane, labelled \textbf{S} (save) or
\textbf{L} (load) with the checkpoint name above the letter. Gates between a
matching S and L still run, but their effect is discarded when the checkpoint
is loaded.

